% ------------------------------------------------------------------------------
\begin{anexo}
\section{Para soporte técnico}
% ------------------------------------------------------------------------------

\subsection{Comprobar que el instalador llegó completo}
\label{sec:huella}

En la página de versiones del repositorio (\href{\enlaceversiones}{\ruta{\enlaceversiones}}) se publica, junto con el instalador, su \textbf{huella SHA-256} en el archivo \ruta{\archivoinstalador.sha256.txt}.
Si la huella calculada en su computadora es igual, el archivo llegó completo y sin cambios.

\begin{enumerate}
	\item Abra la carpeta donde está el instalador.
	\item Haga clic en la barra de dirección de la carpeta, escriba \textbf{powershell} y pulse \tecla{Enter}.
	\item En la ventana que se abre escriba la siguiente línea y pulse \tecla{Enter}:

	\begin{center}
		\ruta{Get-FileHash .\textbackslash \archivoinstalador}
	\end{center}

	\item Compare el texto de la columna \textbf{Hash} con el del archivo \ruta{.sha256.txt}: deben ser iguales.
\end{enumerate}

\subsection{Instalación sin ventanas}

Para instalar en varias computadoras, el instalador acepta los parámetros de Windows Installer. Con \ruta{/qn} instala
sin mostrar el asistente, para el usuario que lo ejecuta y con los dos accesos directos:

\begin{center}
	\ruta{\archivoinstalador /qn}
\end{center}

La desinstalación se hace desde la lista de aplicaciones instaladas (sección «Desinstalar»).

\subsection{Datos técnicos del instalador}

\begin{table}[H]
	\centering
	\caption{Características del instalador.}
	\label{tab:instalador}
	\begin{tabular}{L{4.6cm}L{9.2cm}}
		\hline
		Tipo & Instalador de Windows (MSI dentro de un \ruta{.exe}), creado con jpackage y WiX Toolset. \\
		Alcance & Por usuario; no requiere permisos de administrador. \\
		Carpeta por defecto & \ruta{\%LOCALAPPDATA\%\textbackslash ConcretoLab} \\
		Accesos directos & Menú Inicio (siempre) y escritorio (opcional). \\
		Contenido & La aplicación y su propio entorno Java (no usa el Java instalado en la computadora). \\
		Actualizaciones & Una versión nueva reemplaza a la anterior y conserva los datos. \\
		Firma digital & El instalador no está firmado; por eso Windows puede mostrar el aviso de SmartScreen. \\ \hline
	\end{tabular}
\end{table}

\subsection{Soporte e información legal}

Para informar un problema o sugerir una mejora, use la sección \textit{Issues} del repositorio del proyecto
(\href{\repositorioweb/issues}{\ruta{\repositorioweb/issues}}). Indique la versión instalada (menú \menu{Ayuda}
\textgreater{} \menu{Acerca de ConcretoLab}), la versión de Windows y los pasos que llevan al problema. El registro técnico
del programa está en \ruta{C:\textbackslash Users\textbackslash <su usuario>\textbackslash .concretolab\textbackslash logs}.

ConcretoLab Desktop es una herramienta de apoyo para el anteproyecto y la enseñanza. Sus resultados, en particular los
del análisis sísmico, son preliminares y no están validados para diseñar ni construir obras: deben ser revisados y
firmados por un ingeniero civil colegiado. El instalador incluye componentes de terceros (Java, JavaFX, Spring Boot,
SQLite y otros) distribuidos con sus propias licencias, que se detallan en \menu{Acerca de ConcretoLab}.
\end{anexo}
